\documentclass[11pt]{article}
\usepackage[utf8]{inputenc}
\usepackage{amsmath,amssymb}
\usepackage[margin=1in]{geometry}
\usepackage{hyperref}
\emergencystretch=3em

\title{Two-term asymptotic expansion of the chart standard integral}
\author{Claude, with Daniel Murfet}
\date{2026-09-06}

\begin{document}
\maketitle
\sloppy

\begin{abstract}
The $O$-form of the second-order result: for $\xi(u)=\xi_0+cu^m$,
\[
Z_{\rm chart}(n)=\tfrac1{2k}S_{\mu_1}(\xi_0)\,n^{-\mu_1}+\beta c\,\tfrac1{2k}S_{\mu_2+1/2}(\xi_0)\,n^{-\mu_2}
+O\bigl(n^{-\mu_3}\bigr)\qquad(n\to\infty),
\]
with $\mu_1,\mu_2,\mu_3=(h{+}1)/(2k),\,(h{+}m{+}1)/(2k),\,(h{+}2m{+}1)/(2k)$ three consecutive points of
$\Lambda(h,k)$. The exponentially small boundary tails --- including the $\sqrt n\,e^{-\varepsilon n}$ one
from the first-order insertion --- are absorbed into the $O(n^{-\mu_3})$. Zero \texttt{sorry}/\texttt{axiom}.
\end{abstract}

\section{Setting}
Unit~28 bounds the two-term error by $C_2n^{-\mu_3}+\text{tail}_0+|\beta c|\,\text{tail}_1$ for $n\ge1$,
with the tails explicit integrals. To read this as an asymptotic expansion in the paper's sense one
needs the tails to be $O(n^{-\mu_3})$: $\text{tail}_0\le K_0e^{-\varepsilon n}$ directly from unit~25,
while $\text{tail}_1=\sqrt n\int_b^\infty u^{h+m+k}e^{\rm pop}\le\sqrt n\,K_1e^{-\varepsilon n}$, and both
$e^{-\varepsilon n}$ and $\sqrt n\,e^{-\varepsilon n}$ are $o(n^{-\mu_3})$.

\section{The thing formalised}
{\small
\begin{verbatim}
theorem standardIntegral1D_second_order_isBigO (β ξ₀ c b : ℝ) (h k m : ℕ)
    (hβ : 0 < β) (hb : 0 < b) (hk : 0 < k) :
    (fun n : ℝ => Z_chart n
        - (1/(2k)) n^(-μ₁) S_{μ₁}(ξ₀) - β c ((1/(2k)) n^(-μ₂) S_{μ₂+1/2}(ξ₀)))
      =O[atTop] fun n : ℝ => n ^ (-μ₃)
\end{verbatim}
}

\section{Proof strategy}
\texttt{IsBigO.of\_bound} with the constant $C_2+K_0+|\beta c|K_1$, eventually for $n\ge1$. The two
little-$o$ facts are \texttt{isLittleO\_exp\_neg\_mul\_rpow\_atTop} at exponent $-\mu_3$, and for the
$\sqrt n$-weighted one, \texttt{IsBigO.mul\_isLittleO} of $\sqrt n=O(\sqrt n)$ with the same lemma at
$-\mu_3-\tfrac12$, followed by \texttt{IsLittleO.congr'} using $\sqrt n\cdot n^{-\mu_3-1/2}=n^{-\mu_3}$
for $n>0$ (\texttt{Real.sqrt\_eq\_rpow}, \texttt{rpow\_add}). Their \texttt{.bound one\_pos} instances are
pulled into the same \texttt{filter\_upwards} as $n\ge1$; the rest is the unit-28 inequality with the
tails replaced by $K_ie^{-\varepsilon n}$ (\texttt{standardIntegral1D\_tail\_le} at $h$ and at $h+m+k$)
and then by $K_in^{-\mu_3}$, finished by \texttt{gcongr} and \texttt{ring}.

\section{Roadblocks and resolutions}
None: first-try compile. Idioms worth noting --- \texttt{.bound one\_pos} produces norms of
un-beta-reduced lambda applications, so convert with \texttt{simp only [Real.norm\_eq\_abs]} (which
beta-reduces) before \texttt{rw [abs\_of\_pos \dots]}; the first-order tail integrand is rewritten to
$\sqrt n\cdot u^{h+m+k}e^{\rm pop}$ by \texttt{funext} + \texttt{pow\_add} so that the plain tail lemma
applies; and \texttt{set μ₃} at the top folds the goal but not later-derived hypotheses, so
\texttt{rw [← hμ₃] at hmain} re-aligns them.

\section{What was Mathlib, what was new}
Mathlib: \texttt{IsBigO.of\_bound}, \texttt{IsBigO.mul\_isLittleO}, \texttt{IsLittleO.congr'},
\texttt{IsLittleO.bound}, \texttt{isLittleO\_exp\_neg\_mul\_rpow\_atTop}, \texttt{Real.Gamma\_pos\_of\_pos}.
New: the two-term asymptotic expansion of the chart standard integral, the first formal instance of
\texttt{cor:standardintegralexp} beyond leading order.

\section{Lessons}
Keep the explicit-inequality theorem (unit~28) and the $O$-packaging (this unit) separate: the former
is what one computes with (constants visible), the latter is what one quotes; the packaging step is
short once the little-$o$ scales are named.

\section{Follow-ups}
Higher-order terms follow the same pattern from \texttt{standardIntegral1D\_monomial\_expansion} once the
remainder $\sum_{p\ge N}$ is bounded; the multivariate tree remains the substantive gap.
\end{document}
